% `\scantokens` opens a *file*, and `\everyeof` is what ends it.
%
% e-TeX spells `\scantokens` `pseudo_start` and `etex.ch` puts it in §367's
% `input` case — the one `expand` handles — so the pseudo-file starts *during
% expansion*, and its material becomes part of whatever is being built.
% Dispatched from `main_control` instead it works only at the top level, and
% `\edef\a{\scantokens{ab}}` keeps the literal `\scantokens {ab}`.
%
% The pseudo-file is a file in every other respect too. Each of its lines gets
% `\endlinechar` appended, the last one included (§31 reads lines, §360 appends
% the character), and `\everyeof` goes in when it ends — `etex.ch` in §362's
% `force_eof` branch:
%
%     if every_eof<>null then if not eof_seen[index] then
%       begin begin_token_list(every_eof,every_eof_text);
%       eof_seen[index]:=true; goto restart;
%       end;
%
% Once per file: the *second* arrival at the end is the one that closes it.
%
% expl3's `\tl_set_rescan:Nnn` is all of this at once — it parks a terminator
% in `\everyeof`, wraps the material in markers, and reads it back inside a
% `\tl_set:No`. LaTeX rescans a `\def` for every deferred definition, and
% hyperref installs one on `\@thm`, so without `\everyeof` every theorem in a
% document loading hyperref came out as the literal text of its own
% `\begin{thm}[…]`.
%
% Only forms real TeX accepts without complaint are here. `\edef` around a
% `\scantokens` is *not* one of them on its own — a file may not end in the
% middle of a definition, and TeX says so — which is why expl3 closes the group
% from inside `\everyeof`, and why `\everyeof{\noexpand}` appears below.
\input plain
\catcode`\@=11
% Catcodes are applied afresh, so a `#` that arrived as catcode 12 comes back
% as a parameter character — which is the whole point of the primitive.
\scantokens{\message{<top=read>}}
% `\everyeof` fires when the pseudo-file ends, and only for the file that ended.
\everyeof{\message{<eof=fired>}}
\scantokens{\message{<inner=read>}}
\everyeof{}
\scantokens{\message{<after=read>}}
% …and once per file. A pseudo-file with two lines still ends once.
\newlinechar=`\^^J
\everyeof{\message{<eof2=fired>}}
\scantokens{\message{<line1>}^^J\message{<line2>}}
\everyeof{}
\newlinechar=-1
% Expandable: §367 handles `\scantokens` in `expand`, so the pseudo-file starts
% while the `\edef` body is being built. `\everyeof{\noexpand}` is what makes
% that legal — a file may not end in the middle of a definition, and the
% `\noexpand` reads the `}` that follows and backs it up, so `scan_toks` sees a
% brace and closes the body (§369 attaches the `\notexpanded:` marker only to a
% control sequence).
\endlinechar=-1
\everyeof{\noexpand}
\edef\a{\scantokens{ab}}\message{<edef=\meaning\a>}
% …and the pseudo-file's last line still gets `\endlinechar`, which is why the
% first of these ends in a space and the second does not.
\endlinechar=13
\edef\b{\scantokens{ab}}\message{<endline=\meaning\b>}
\everyeof{}
% One step of expansion has to start the file too, or a delimited argument
% scanned afterwards swallows the literal `\scantokens {…}` — which is what
% expl3's `\tl_set_rescan:Nnn` does, and it reads the rest of the document.
\endlinechar=-1
\everyeof{\noexpand}
% The delimiter goes *inside* the pseudo-file: a delimited argument may not
% span the end of one, and `\everyeof{\noexpand}` does not help there.
\def\grab#1STOP{\message{<grabbed=\detokenize{#1}>}}
\expandafter\grab\scantokens{abSTOP}
\everyeof{}
\endlinechar=13
\end
